\documentclass[../../main.tex]{subfiles}
\begin{document}

\begin{flushright}
\footnotesize\textit{【自動文字起こし・要確認】}
\end{flushright}

[解] (1) はじめの操作で場合分けする。
$1^\circ \ 0 \le x \le 15$
$P_m(x) = \dfrac{1}{2} P_{m-1}(2x) + \dfrac{1}{2} P_{m-1}(0)$ であり、$P_{m-1}(0) = 0$ だから、
\[ P_m(x) = \dfrac{1}{2} P_{m-1}(2x) \]

$2^\circ \ 16 \le x \le 30$
$P_m(x) = \dfrac{1}{2} P_{m-1}(2x-30) + \dfrac{1}{2} P_{m-1}(30)$ であり、$P_{m-1}(30) = 1$ から、
\[ P_m(x) = \dfrac{1}{2} P_{m-1}(2x-30) + \dfrac{1}{2} \]

以上まとめて、
\[ P_m(x) = \begin{cases} \dfrac{1}{2} P_{m-1}(2x) & (0 \le x \le 15) \\ \dfrac{1}{2} P_{m-1}(2x-30) + \dfrac{1}{2} & (16 \le x \le 30) \end{cases} \]

(2) (1)から、
\[ P_{2n+2}(10) = \dfrac{1}{2} P_{2n+1}(20) = \dfrac{1}{2}\left(\dfrac{1}{2} P_{2n}(10) + \dfrac{1}{2}\right) = \dfrac{1}{4} P_{2n}(10) + \dfrac{1}{4} \]
であり、
\[ P_{2n+2}(10) - \dfrac{1}{3} = \dfrac{1}{4} \left(P_{2n}(10) - \dfrac{1}{3}\right) \]
となる。$P_2(10) = \dfrac{1}{4}$ だから、くり返し用いて、
\[ P_{2n}(10) = \left(\dfrac{1}{4}\right)^{n-1} \left(\dfrac{1}{4} - \dfrac{1}{3}\right) + \dfrac{1}{3} = \dfrac{1}{3} \left\{ 1 - \left(\dfrac{1}{4}\right)^n \right\} \]

(3) (1) から
\begin{align}
P_{4n+4}(6) &= \dfrac{1}{2} P_{4n+3}(12) = \dfrac{1}{4} P_{4n+2}(24) = \dfrac{1}{4}\left(\dfrac{1}{2} P_{4n+1}(18) + \dfrac{1}{2}\right) \\
&= \dfrac{1}{8} \left(\dfrac{1}{2} P_{4n}(6) + \dfrac{1}{2}\right) + \dfrac{1}{8} = \dfrac{1}{16} P_{4n}(6) + \dfrac{3}{16}
\end{align}
\[ \therefore P_{4n+4}(6) - \dfrac{1}{5} = \dfrac{1}{16} \left(P_{4n}(6) - \dfrac{1}{5}\right) \]
となる。$P_4(6) = \dfrac{3}{16}$ だから、くり返し用いて、
\[ P_{4n}(6) = \left(\dfrac{1}{16}\right)^{n-1} \left(\dfrac{3}{16} - \dfrac{1}{5}\right) + \dfrac{1}{5} = \dfrac{1}{5} \left\{ 1 - \left(\dfrac{1}{16}\right)^n \right\} \]

\end{document}
